\chapter{概率与抽样}\label{ch:g10:proba}

概率为偶然性赋予数值。从
具有等可能结果的试验（掷骰子、掷硬币、抽扑克牌）出发，本章建立
计算事件概率的基本法则，引入两步试验的树状图，
并以概率与观测频率之间的联系收尾。
本章理论后续将在
\cref{ch:g11:prob}中展开。

\section{随机试验与事件}

\begin{definition}[结果与事件]\label{def:g10:proba:events}
称\emph{随机试验}\index{随机试验}为结果无法预先确定的试验
（如掷一颗骰子）。所有可能结果组成的集合称为
\emph{样本空间}\index{样本空间}
$\Omega = \{\omega_1, \dots, \omega_n\}$；其元素称为
\emph{结果}。称\emph{事件}\index{事件}为由若干结果组成的集合；
当试验结果属于该集合时，事件\emph{发生}。
\end{definition}

\begin{definition}[概率分布]\label{def:g10:proba:distribution}
对每个结果 $\omega_i$ 赋予一个数 $p_i \geq 0$，使得
$p_1 + \dots + p_n = 1$，即定义了一个\emph{概率}\index{概率}：
$\P(A)$ 即事件 $A$ 的概率，等于各 $p_i$ 之和，
其中 $A$ 中的结果。特别地，$\P(\Omega) = 1$，不可能事件
$\varnothing$ 的概率为 $0$。
\end{definition}

\begin{proposition}[等可能结果]\label{prop:g10:proba:uniform}
若 $n$ 个结果等可能（每个 $p_i = \frac1n$），则对任一
事件 $A$，有：
\[
\P(A) = \frac{\text{ 结果数 } A}{n}
= \frac{\text{ 有利情形数 }}{\text{ 可能情形总数 }} .
\]
\end{proposition}

\begin{proof}
$A$ 含有 $k$ 个结果，每个结果的概率均为 $\frac1n$，
故 $\P(A) = \frac kn$。
\end{proof}

\begin{example}\label{ex:g10:proba:die}
掷一枚均匀骰子：$\Omega = \{1, 2, 3, 4, 5, 6\}$，每个结果的
概率均为 $\frac16$。事件 $A$：“点数为偶数”，则
$A = \{2, 4, 6\}$，于是 $\P(A) = \frac36 = \frac12$。事件 $B$：“点数
不小于 $5$”，则 $B = \{5, 6\}$，于是 $\P(B) = \frac26 =
\frac13$。
\end{example}

\section{事件的运算}

\begin{definition}[交集、并集、对立事件]\label{def:g10:proba:operations}
设 $A$、$B$ 为事件。
\begin{itemize}
  \item $A \cap B$（“$A$ 且 $B$”）表示两者同时发生；
  \item $A \cup B$（“$A$ 或 $B$”）表示二者至少有一个发生；
  \item \emph{对立事件}\index{对立事件} $\bar A$（“非 $A$”），
        当且仅当 $A$ 不发生时发生；
  \item $A$ 与 $B$ 称为\emph{互斥}\index{互斥}，
        即二者不能同时发生：$A \cap B = \varnothing$。
\end{itemize}
\end{definition}

\begin{theorem}[加法公式]\label{thm:g10:proba:addition}
对任意事件 $A$、$B$，有：
\[
\P(A \cup B) = \P(A) + \P(B) - \P(A \cap B),
\qquad
\P(\bar A) = 1 - \P(A) .
\]
若事件互斥，则第一式化简为
$\P(A \cup B) = \P(A) + \P(B)$。
\end{theorem}

\begin{proof}
将 $\P(A) + \P(B)$ 相加时，$A \cup B$ 中每个结果的概率
恰被计入一次，\emph{除} $A \cap B$ 中的结果外
这些结果被计入两次——一次计入 $A$，一次计入 $B$。减去 $\P(A \cap B)$ 即可消除重复计数。
对于对立事件，$A$ 与 $\bar A$ 互斥，
且 $A \cup \bar A = \Omega$，故 $\P(A) + \P(\bar A) = \P(\Omega) = 1$。
\end{proof}

\begin{omfigure}
\begin{tikzpicture}[scale=0.9]
  \draw[omExo, thick, rounded corners] (-3.1,-1.7) rectangle (3.4,1.7);
  \node[omExo, font=\small] at (3.0,1.45) {$\Omega$};
  \begin{scope}
    \clip (-0.8,0) circle (1.25);
    \fill[omProp!25] (0.8,0) circle (1.25);
  \end{scope}
  \draw[omDef, thick] (-0.8,0) circle (1.25);
  \draw[omThm, thick] (0.8,0) circle (1.25);
  \node[omDef, font=\small] at (-1.6,0.9) {$A$};
  \node[omThm, font=\small] at (1.7,0.9) {$B$};
  \node[font=\small] at (0,0) {$A \cap B$};
\end{tikzpicture}

{\small 两个相交事件：在 $\P(A) + \P(B)$ 中，阴影所示的
交集 $A \cap B$ 被计入两次，这
正是加法公式中要减去它的原因。}
\end{omfigure}

\begin{example}\label{ex:g10:proba:cards}
从一副标准的 $52$ 张扑克牌中抽取一张。设 $A$：“抽到的牌是
红心”（$13$ 张），$B$：“抽到的牌是国王”（$4$ 张）。则
$A \cap B$ 为“红心国王”（$1$ 张），且
\[
\P(A \cup B) = \frac{13}{52} + \frac{4}{52} - \frac{1}{52} = \frac{16}{52}
= \frac{4}{13}.
\]
用对立事件公式往往最快：抽到的牌
\emph{不是}红心的概率为 $1 - \frac{13}{52} = \frac34$。
\end{example}

\section{树状图}

\begin{method}[两步实验]\label{met:g10:proba:tree}
对于分两步进行的实验：
\begin{enumerate}
  \item 画一个\emph{树状图}\index{树状图}：第一步的每个可能结果对应一条分支，
        然后从每个结果出发，第二步的每个可能结果各对应一条分支，
        并在每条分支上标明概率
        （从任一节点发出的各分支概率之和为 $1$）；
  \item \emph{路径}（树状图的叶节点）的概率，
        等于沿其各分支概率的乘积；
  \item 事件的概率等于使其发生的各条
        路径的概率之和。
\end{enumerate}
\end{method}

\begin{example}\label{ex:g10:proba:tree}
一个罐子中装有 $3$ 个红球和 $2$ 个蓝球。取出一球后将其放回，
再取一次。每次取到红球（R）的概率为 $\frac35$，
取到蓝球（B）的概率为 $\frac25$。

\begin{omfigure}
\begin{tikzpicture}[scale=1.0]
  \coordinate (root) at (0,0);
  \coordinate (R) at (2.3,1.1);
  \coordinate (B) at (2.3,-1.1);
  \coordinate (RR) at (5.0,1.7);
  \coordinate (RB) at (5.0,0.5);
  \coordinate (BR) at (5.0,-0.5);
  \coordinate (BB) at (5.0,-1.7);
  \draw (root) -- (R) node[midway, above, font=\small, sloped] {$\frac35$};
  \draw (root) -- (B) node[midway, below, font=\small, sloped] {$\frac25$};
  \draw (R) -- (RR) node[midway, above, font=\small, sloped] {$\frac35$};
  \draw (R) -- (RB) node[midway, below, font=\small, sloped] {$\frac25$};
  \draw (B) -- (BR) node[midway, above, font=\small, sloped] {$\frac35$};
  \draw (B) -- (BB) node[midway, below, font=\small, sloped] {$\frac25$};
  \node[omThm, font=\small] at (R) [above left=-2pt] {R};
  \node[omDef, font=\small] at (B) [below left=-2pt] {B};
  \node[right, font=\small] at (RR)
    {RR:\ $\frac35 \times \frac35 = \frac{9}{25}$};
  \node[right, font=\small] at (RB)
    {RB:\ $\frac35 \times \frac25 = \frac{6}{25}$};
  \node[right, font=\small] at (BR)
    {BR:\ $\frac25 \times \frac35 = \frac{6}{25}$};
  \node[right, font=\small] at (BB)
    {BB:\ $\frac25 \times \frac25 = \frac{4}{25}$};
\end{tikzpicture}

{\small 有放回取两次的树状图。四条路径
概率之和为 $1$。}
\end{omfigure}

两球颜色相同的概率为
$\P(\text{RR}) + \P(\text{BB}) = \frac{9}{25} + \frac{4}{25} =
\frac{13}{25}$；至少有一个蓝球的概率为
$1 - \P(\text{RR}) = 1 - \frac{9}{25} = \frac{16}{25}$（对立事件公式）。
\end{example}

\begin{remark}
不放回取球时，第二次取球的概率依赖于第一次的结果：
在上述罐子中，若第一次取到红球，则只剩 $2$ 个红球和 $2$ 个
蓝球，因此第二层分支上的两个概率分别为 $\frac24$ 和 $\frac24$。
树状图方法仍然适用，只是第二层分支上的数值
发生变化。
\end{remark}

\section{频率与抽样}

在随机试验重复进行多次时，某事件的频率
趋近于该事件的概率 --- 这正是概率能够
描述现实世界的原因。

\begin{definition}[样本]\label{def:g10:proba:sample}
容量为 $n$ 的\emph{样本}\index{样本}是指将同一试验
独立重复进行 $n$ 次所得的结果。样本中某事件的\emph{频率}
是该事件出现的次数除以 $n$。
\end{definition}

\begin{example}\label{ex:g10:proba:sampling}
一枚均匀硬币正面朝上的概率为 $0.5$。将它掷 $100$ 次，
很少恰好得到 $50$ 次正面：样本会\emph{波动}。典型的
$100$ 次掷币得到的频率如 $0.46$、$0.53$、$0.49$ --- 接近
$0.5$，但并不相等。样本越大，波动越小：掷 $10\,000$ 次时，
观察频率通常以 $0.01$ 以内的偏差落在 $0.5$ 附近。
\end{example}

\begin{remark}[波动区间]
一条实用的经验法则（在\cref{ch:g11:binom}中证明）：
对于容量为 $n$ 的样本，概率为 $p$ 的事件（$n$ 较大时），
观察频率落在下列区间内
\[
\intcc{p - \frac{1}{\sqrt n}}{\,p + \frac{1}{\sqrt n}}
\]
约占 $95\%$ 的样本。若某个观察频率远远落在
该区间之外，则应怀疑 $p$ 的取值，这正是
统计检验的出发点。
\end{remark}

\begin{example}\label{ex:g10:proba:test}
某工厂声称其零件中只有 $10\%$ 为次品
（$p = 0.1$）。在容量 $n = 400$ 的样本中，有 $64$ 件次品，
观察频率为 $\frac{64}{400} = 0.16$。波动区间为
$\intcc{0.1 - \frac{1}{20}}{0.1 + \frac{1}{20}} = \intcc{0.05}{0.15}$,
$0.16$ 落在该区间之外，该样本对这一说法提出了严重质疑。
\end{example}

\section{练习}

\begin{exercise}[$\star$]\label{exo:g10:proba:1}
掷一枚均匀的骰子。求下列各事件的概率：“掷出
$6$”；“掷出奇数”；“点数不超过 $4$”；“掷出
$7$”。
\end{exercise}

\begin{exercise}[$\star$]\label{exo:g10:proba:2}
袋中装有 $5$ 颗绿色、$3$ 颗黄色和 $2$ 颗黑色的弹珠，随机取出一颗。
求取出的一颗弹珠是绿色的概率、不是黑色的概率，
以及是绿色或黄色的概率。
\end{exercise}

\begin{exercise}[$\star$]\label{exo:g10:proba:3}
某班有 $30$ 名学生，其中 $18$ 人学习西班牙语，$10$ 人学习德语，
$4$ 人两门都学。随机抽取一名学生，利用加法公式，
求该学生至少学习两门语言中的一门的概率，
再求该学生两门语言都不学的概率。
\end{exercise}

\begin{exercise}[$\star$]\label{exo:g10:proba:4}
掷一枚均匀的硬币两次。画出该试验的树状图，并求
两次都掷出正面的概率；恰有一次正面的概率；
至少有一次正面的概率。
\end{exercise}

\begin{exercise}[$\star\star$]\label{exo:g10:proba:5}
袋中有 $4$ 个红球和 $6$ 个白球。依次取出两个球，
且每次取出后\emph{不放回}。画出树状图，在第二层标出正确的概率，
并求取出两个红球的概率，
以及取出两个颜色不同的球的概率。
\end{exercise}

\begin{exercise}[$\star\star$]\label{exo:g10:proba:6}
掷两枚均匀的骰子，并记录两个点数之和。
\begin{enumerate}
  \item 说明为何样本空间可取为 $36$ 个等可能的
        数对 $(1,1), (1,2), \dots, (6,6)$。
  \item 求点数之和为 $7$ 的概率；为 $12$ 的概率；
        不超过 $4$ 的概率。
  \item 哪一个点数之和最有可能出现？
\end{enumerate}
\end{exercise}

\begin{exercise}[$\star\star$]\label{exo:g10:proba:7}
某道选择题给出 $4$ 个选项，其中只有一个正确。某学生
完全随机地作答 $2$ 道这样的独立选择题。求两题全部答对、
恰答对一题、两题全部答错的概率；
并验证这三个概率之和为 $1$。
\end{exercise}

\begin{exercise}[$\star\star$]\label{exo:g10:proba:8}
某政客宣称支持率为 $50\%$。在对 $n = 900$ 名随机选取的选民所做的民意调查中，
有 $396$ 人表示支持。
\begin{enumerate}
  \item 求观测频率，以及
        围绕 $p = 0.5$，在 $n = 900$ 时的波动区间。
  \item 该民意调查是否对上述宣称提出了质疑？
\end{enumerate}
\end{exercise}

\begin{exercise}[$\star\star\star$]\label{exo:g10:proba:9}
玩一次游戏需付 $2$：掷一枚均匀的骰子，
赢得 $6$ 需掷出 $6$，赢得 $3$ 需掷出 $5$，否则一无所得。
\begin{enumerate}
  \item 求各赢得金额（$6$、$3$、$0$）的概率。
  \item 进行 $600$ 次游戏，你预计每种结果大约各出现多少次？
        求总的期望收益与损失：该游戏
        是否值得参加？
\end{enumerate}
\end{exercise}

\section{周末问题：骗过一千位数学家的
三扇门}

\begin{problem}[{周末问题 --- 蒙提霍尔悖论及其同类：信息如何改变概率，
以及样本如何
将主张与一厢情愿的想法区分开来}]
\label{pb:g10:proba:1}
1990 年，专栏作家玛丽莲·沃斯·莎凡特发表了一道谜题，内容涉及
一位游戏节目主持人、三扇门、一辆汽车和两只山羊 --- 她给出了正确答案。
一万名读者（其中包括数百位数学博士）纷纷来信，
坚称她的答案是错的。本题将借助树状图
（\cref{met:g10:proba:tree}）、对立事件及加法公式
（\cref{thm:g10:proba:addition}），
让你亲自判定这场争论，
并以实践中裁决此类争议的工具作结：
波动的样本。

\medskip
\noindent\textbf{第一部分 --- 规则热身。}
\begin{enumerate}
  \item 从一副标准的 $52$ 张扑克牌中抽取一张，求
        $P(\text{国王或红心})$ --- 用加法公式计算，
        并说明为何仅凭 $\frac{4}{52} + \frac{13}{52}$
        是错误的。
  \item 抛掷三次均匀的硬币，利用对立事件求
        $P(\text{至少出现一次正面})$。
  \item 袋中有 $3$ 颗红色弹珠和 $2$ 颗蓝色弹珠，从中
        不放回地取出两颗。借助树状图，求
        $P(\text{两颗红色})$ 和 $P(\text{一红一蓝})$。
  \item 补全第 3 题的概率分布，
        求 $P(\text{两颗蓝色})$，并验证三个概率
        之和为 $1$。
  \item 有人报告，对于两个 \emph{互斥}事件，
        $P(A) = 0.6$，$P(B) = 0.5$。
        试用一行运算指出其结论有误。
\end{enumerate}

\medskip
\noindent\textbf{第二部分 --- 三扇门。}
游戏规则：汽车藏在三扇门之一的后面，山羊藏在
另外两扇门后。你选定一扇门（不妨设为 1 号门）。主持人 --- 他
\emph{知道}汽车的位置 --- 打开其余两扇门中的一扇，
门后总是露出一只山羊，并给你机会
换选剩下的那扇未打开的门。
\begin{enumerate}[resume]
  \item 在任何计算之前：换选有利、有害，还是
        无关紧要？写下你的直觉答案 --- 请诚实作答。
  \item 现在画出树状图，按汽车的真实位置分类讨论
        （各情况的概率均为 $\frac13$）：若汽车在你所选的门后，
        换选会输；若汽车在另外两扇门中的任一扇后，
        主持人须怎样做？此时换选能否获胜？
        求 $P(\text{换选获胜})$。
  \item 一句话概括：换选恰好在你第一次的选择
        有误时获胜。请再次得出结论。
  \item 现在看关键的变体 --- “蒙提失误”：主持人
        并\emph{不}知道汽车的位置，随机打开另外两扇门中的一扇，
        碰巧露出一只山羊。列出六种等可能的情形
        （汽车位置 $\times$ 打开的门），
        剔除汽车被露出的情形，
        并求在\emph{剩余的情形中}换选获胜的概率。
        什么发生了改变，
        以及主持人的知情与否为何重要？
  \item 直觉启发：一百扇门，你选定其中一扇，知情的
        主持人打开 $98$ 扇有山羊的门。是否换选？换选
        获胜的概率是多少？
  \item 你面对一位怀疑者，他与那一千位写信者一样，
        坚持认为是“五五开”。请描述一个基于纸牌的模拟
        （三张牌、多轮、由一位朋友担任主持人），并说明
        换选获胜的比例预计为多少，
        在 $300$ 轮之后 --- 以本章的波动区间作为
        你的误差范围。
\end{enumerate}

\medskip
\noindent\textbf{第三部分 --- 悖论的近亲。}
\begin{enumerate}[resume]
  \item 某家庭有两个孩子，已知\emph{至少有
        一个是男孩}。列出与该信息相容的等可能两孩情形，
        并计算
        $P(\text{两个男孩})$。
  \item 同一家庭，但现在已知：\emph{年长的孩子
        是男孩}。重新计算。两个答案为何不同？它们与
        两个蒙蒂有何共同之处？
  \item 三枚均匀硬币：计算 $P(\text{三枚结果全相同})$。
        数学家达朗贝尔于 1754 年主张另一个
        答案：“要么三枚全相同，要么不全相同——
        两种情形中居其一。”请指出确切的错误，并结合
        \cref{prop:g10:proba:uniform} 加以说明。
  \item 一次游戏收费 $2$ 欧元：掷一枚骰子，掷出六点赢 $6$ 欧元，
        掷出五点赢 $3$ 欧元，其余情形一无所得。
        （见\cref{exo:g10:proba:9}）。若进行 $600$ 局游戏，
        计数预测的每局平均收益是多少？一位保险
        精算师用的也是同样的数学：每年发生概率为 $1\,\%$ 的
        一笔 $10\,000$ 欧元理赔。其
        \emph{公平}保费为多少？实际保费为何
        高于它？
\end{enumerate}

\medskip
\noindent\textbf{第四部分 --- 以样本为裁判。}
\begin{enumerate}[resume]
  \item 掷一枚硬币 $100$ 次。利用波动区间
        $p \pm \frac{1}{\sqrt n}$，判断下列结果中哪些
        值得怀疑：$55$ 次正面朝上；$70$
        次正面朝上。
  \item 一项对 $1\,000$ 名选民的民调显示，某候选人得票率为 $52\,\%$。
        计算以 $50\,\%$ 为中心的波动区间，并
        说明该民调是否\emph{证明}该候选人
        领先。
  \item 某工厂声称其零件次品率为 $4\,\%$。抽取一个
        容量为 $400$ 的样本，其中有 $28$ 件次品（$7\,\%$）；之后
        又抽取容量为 $2\,500$ 的样本，其中有 $175$ 件（仍为 $7\,\%$）。分别用两个样本检验该
        说法。两次结论为何不同？
  \item 每项调查都受两类误差困扰：抽样
        \emph{偏差}（初中分册中的周末问题“数据如何说谎”）和抽样
        \emph{波动}（本章内容）。随着
        $n$ 增大，哪一种误差会缩小？哪一种
        无论样本多大都不会消失？各用一句话回答。
  \item 总结 --- 概率学者的检查清单，每项一行：
        计数之前，先确认结果等可能（第 15 题）；
        多步骤问题与“至少”类问题，用树状图与对立事件
        （第 2--3 题）；\emph{信息改变概率，获得
        信息的方式至关重要}（第 7--14 题）；
        样本的波动大约为
        $\frac{1}{\sqrt n}$，由模拟来裁决
        （第 11、16--18 题）。
\end{enumerate}
\end{problem}
